\begin{table}[htbp]
\centering
\caption{Border premium including ZIP codes with no dispensary activity}
\label{tab:nm_baseline_zeroinclusive}

\begin{tabular*}{\textwidth}{@{\extracolsep{\fill}}lccc}
\toprule
& (1) & (2) & (3) \\
& \shortstack{Revenue\\per capita}
& \shortstack{Dispensary\\density}
& \shortstack{Revenue\\per dispensary} \\
& $\operatorname{asinh}(y)$
& $\operatorname{asinh}(y)$
& $\log(y)$ \\
\midrule

Border ZIP
    & $0.874^{***}$
    & $0.257^{***}$
    & $0.255^{*}$ \\
    & (0.226)
    & (0.075)
    & (0.149) \\

\addlinespace

$\log(\text{Population})$
    & $0.448^{***}$
    & 0.012
    & $0.242^{***}$ \\
    & (0.083)
    & (0.020)
    & (0.051) \\

\midrule

Quarter fixed effects
    & Yes & Yes & Yes \\

Observations
    & 6,624
    & 6,624
    & 1,704 \\

Number of ZIP codes
    & 368
    & 368
    & 116 \\

$R^{2}$
    & 0.146
    & 0.034
    & 0.167 \\

\bottomrule
\end{tabular*}

\begin{minipage}{\textwidth}
\vspace{0.4em}
\fontsize{9pt}{9pt}\selectfont

\textit{Notes:} Companion to Table~\ref{tab:nm_baseline}. Columns~(1) and~(2)
apply the inverse hyperbolic sine transformation to revenue per capita and to
dispensaries per 1{,}000 residents, so that the 368 populated ZIP codes with no
recorded dispensary activity remain in the sample; the border premium is close
to the log-level estimates in Table~\ref{tab:nm_baseline}. Column~(3) is
identical to Table~\ref{tab:nm_baseline}. $\text{Population}$ is the ZIP code's
own resident population. Quarter fixed effects throughout; standard errors,
in parentheses, are clustered at the county level.
\end{minipage}

\end{table}
